% swiftui-state.tex — some View, @State/@Binding/@StateObject/@ObservedObject/@EnvironmentObject,
% iOS 17 Observation, @ViewBuilder, view identity and when body re-runs.
% Sources (checked 2026-10-05):
%   developer.apple.com/documentation/swiftui/state · .../binding · .../stateobject (iOS 14; autoclosure,
%     once per identity) · .../observedobject · .../environmentobject (iOS 13)
%   developer.apple.com/documentation/swiftui/migrating-from-the-observable-object-protocol-to-the-observable-macro
%     (iOS 17: @State, @Bindable, @Environment(T.self), .environment(o))
%   developer.apple.com/documentation/swiftui/viewbuilder · .../view/task(priority:_:) (iOS 15)
%   developer.apple.com/videos/play/wwdc2021/10022 "Demystify SwiftUI" (identity, lifetime, dependencies)
% @labels: area=ios-swift kind=api platform=apple topic=ui,data discussion=70
% @tags: state, binding, stateobject, observedobject, environmentobject, observableobject, published, observable, bindable, viewbuilder, opaque-types, view-identity
\documentclass[8pt]{extarticle}
\usepackage{printup-sheet}

\lstdefinelanguage{SwiftSheet}{
  morekeywords={func,let,var,struct,class,final,enum,protocol,extension,return,if,else,
    guard,try,await,async,throws,case,switch,self,some,any,where,init,in,private,
    @State,@Binding,@StateObject,@ObservedObject,@Published,@Observable,@Bindable},
  alsoletter={@}, sensitive=true, morecomment=[l]{//}, morestring=[b]"}

\begin{document}

\sheettitle{SwiftUI state \& data flow}{ios-swift · folio}

\oneliner{A view is a cheap \textbf{struct} rebuilt constantly; its state lives \textbf{outside} it in
SwiftUI's storage, keyed by \textbf{view identity}. The only question each wrapper answers is
\textbf{who owns the truth}: this view (\texttt{@State}, \texttt{@StateObject}) or someone else
(\texttt{@Binding}, \texttt{@ObservedObject}, \texttt{@EnvironmentObject}).}

\begin{multicols}{2}
\raggedright

\section{How it works}
\begin{itemize}
  \item \textbf{\texttt{View}} is a protocol: \texttt{associatedtype Body: View} +
        \texttt{@ViewBuilder var body: Self.Body \{ get \}}. \texttt{some View} is an
        \textbf{opaque result type}: \emph{one} concrete type (e.g.
        \texttt{VStack<TupleView<(Text, Button<Text>)>>}), chosen by the \emph{implementation},
        known to the compiler, hidden from the reader. Not a box, not dynamic — unlike \texttt{AnyView},
        which erases the type and hurts diffing.
  \item \textbf{\texttt{@State}} — view-\textbf{owned} source of truth for a \emph{value}; storage lives
        in SwiftUI, survives \texttt{body} re-runs; the initial value is used \emph{once}. Mark it
        \texttt{private}. \texttt{\$x} (the \texttt{projectedValue}) is a \texttt{Binding}.
  \item \textbf{\texttt{@Binding}} — \textbf{not} storage: a two-way get/set \emph{reference} to truth
        owned elsewhere. Child writes $\to$ parent's \texttt{@State} changes $\to$ both re-render.
        Passed as \texttt{\$name}; \texttt{.constant(x)} in previews.
  \item \textbf{\texttt{@StateObject}} — view \textbf{creates and owns} an \texttt{ObservableObject};
        the \texttt{@autoclosure} initialiser runs \textbf{once per identity}, the instance survives
        struct re-creation.
  \item \textbf{\texttt{@ObservedObject}} — object is \textbf{injected}; view only subscribes, does
        \emph{not} control lifetime. Initialised inline, it is re-created on every parent redraw.
  \item \textbf{\texttt{@EnvironmentObject}} — injected by an ancestor's
        \texttt{.environmentObject(obj)}, looked up by \emph{type} at run time; missing = crash.
        \texttt{@Environment(\textbackslash.colorScheme)} reads \emph{values} (\texttt{\textbackslash.dismiss}, \texttt{\textbackslash.scenePhase}).
  \item \textbf{Notification}: \texttt{@Published} fires \texttt{objectWillChange} \emph{before} the set
        (whole object). \textbf{iOS 17 \texttt{@Observable}} (macro, Observation) tracks \textbf{each
        property \texttt{body} read} — see the table.
  \item \textbf{\texttt{@ViewBuilder}} — a \emph{result builder}: \texttt{buildBlock} packs children
        into \texttt{TupleView}; \texttt{if/else} $\to$ \texttt{\_ConditionalContent}
        (\texttt{buildEither}); \texttt{if} alone $\to$ \texttt{Optional} (\texttt{buildOptional}).
        That is why \texttt{body} may branch but a plain helper \texttt{func -> some View} may not
        (all returns must be one type) — mark it \texttt{@ViewBuilder}.
  \item \textbf{Identity}: \emph{structural} (position in the tree, which \texttt{if} branch) or
        \emph{explicit} (\texttt{.id(\_:)}, \texttt{ForEach} ids). New identity = state destroyed,
        \texttt{onAppear} again, transition played.
  \item \textbf{\texttt{body} re-runs} when something it \emph{read} changes: a \texttt{@State} /
        \texttt{@Binding} value, an observed object publishes (or an \texttt{@Observable} property it read),
        an \texttt{@Environment} value, or the parent passes different inputs. Then SwiftUI diffs.
\end{itemize}

\section{Example}
\begin{lstlisting}[language=SwiftSheet]
final class VM: ObservableObject { @Published var n = 0 }
struct Screen: View {
  @StateObject private var vm = VM()   // made ONCE, owned
  @State private var on = false        // value, owned
  var body: some View {                // one concrete type
    VStack { Toggle("On", isOn: $on)   // $on : Binding<Bool>
             Row(vm: vm, on: $on) }    // hand truth down
  }
}
struct Row: View {
  @ObservedObject var vm: VM           // injected, NOT owned
  @Binding var on: Bool                // someone else's @State
  var body: some View { Button("+") { vm.n += 1; on = true } }
}
\end{lstlisting}

\columnbreak

\section{Structs die, storage stays}
\begin{tikzpicture}[sheet]
  \tikzset{vw/.style={box, text width=27mm, align=left, font=\footnotesize},
           st/.style={cell, font=\ttfamily\footnotesize, minimum width=27mm, minimum height=6mm,
                      align=center, fill=sheetGreen!10, draw=sheetGreen}}
  % left lane: view structs
  \node[vw] (P) at (0,0) {\textbf{Screen} (struct)\\\texttt{@State on}\\\texttt{@StateObject vm}};
  \node[vw] (C) at (0,-1.75) {\textbf{Row} (struct)\\\texttt{@Binding on}\\\texttt{@ObservedObject vm}};
  \node[vw, draw=sheetRed, fill=sheetRed!6] (B) at (0,-3.45)
       {\textbf{Bad} (struct)\\\texttt{@ObservedObject}\\\texttt{var vm = VM()}};
  % right lane: SwiftUI storage
  \node[st] (S) at (4.6,0) {on = true};
  \node[st, fill=sheetBlue!8, draw=sheetBlue] (O) at (4.6,-1.75) {VM \#1\\[-2pt]{\scriptsize created once}};
  \node[st, fill=sheetRed!6, draw=sheetRed, dashed] (X) at (4.6,-3.45)
       {VM \#2, \#3, \#4 …\\[-2pt]{\scriptsize new per redraw}};
  \begin{scope}[on background layer]
    \node[draw=sheetGrey, dashed, rounded corners=3pt, fill=black!3, inner sep=4pt,
          fit=(S)(O)(X), label={[note, align=center]above:SwiftUI storage\\[-2pt]lives as long as the \emph{identity}}] {};
    \node[draw=none, inner sep=3pt, fit=(P)(C)(B),
          label={[note, align=center]above:view structs\\[-2pt]rebuilt on every parent \texttt{body}}] {};
  \end{scope}
  \draw[->, thick, sheetBlue] ([yshift=2mm]P.east) -- node[above, note, pos=0.5]{owns (\texttt{@State})} ([yshift=2mm]S.west);
  \draw[->, thick, sheetBlue] ([yshift=-2mm]P.east) -- node[above, sloped, note, pos=0.22]{owns} ([yshift=1mm]O.west);
  \draw[<->, thick, sheetOrange] ([yshift=2mm]C.east) -- node[below, sloped, note, text=sheetOrange, pos=0.62]{\texttt{\$on} r/w} ([yshift=-2mm]S.west);
  \draw[->, thick, sheetGrey, dashed] ([yshift=-2mm]C.east) -- node[below, note, pos=0.42]{watches only} ([yshift=-2mm]O.west);
  \draw[->, thick, sheetRed] (B.east) -- node[above, note, text=sheetRed]{re-inits!} (X.west);
\end{tikzpicture}

\vspace{2pt}
{\footnotesize
\begin{tabular}{@{}p{17mm}p{27mm}p{29mm}@{}}
\toprule
\textbf{role} & \textbf{\texttt{ObservableObject}} & \textbf{iOS 17 \texttt{@Observable}} \\
\midrule
owns it     & \texttt{@StateObject}          & \texttt{@State} \\
passed in   & \texttt{@ObservedObject}       & plain \texttt{let} / \texttt{var} \\
needs \texttt{\$x.prop} & \texttt{@ObservedObject} & \texttt{@Bindable} \\
injected    & \texttt{@EnvironmentObject}    & \texttt{@Environment(T.self)} \\
provide     & \texttt{.environmentObject(o)} & \texttt{.environment(o)} \\
publishes   & \texttt{@Published} props      & every stored prop (macro) \\
re-renders  & every observer, any change     & only views that \emph{read} it \\
\bottomrule
\end{tabular}\par}

\section{Traps}
\begin{itemize}
  \trap{\textbf{\texttt{@ObservedObject var vm = VM()}} — created but not owned: each parent redraw
        allocates a fresh VM, state resets, \emph{no warning}. Creator: \texttt{@StateObject}.}
  \trap{\textbf{``\texttt{some View} = any view / like \texttt{AnyView}''} — no: \texttt{some} is
        \emph{one fixed type}; returning \texttt{Text} in one branch and \texttt{Image} in another
        fails without \texttt{@ViewBuilder}.}
  \trap{\textbf{``\texttt{@Binding} is a copy''} — no: it stores nothing; writes go to the owner.
        Two \texttt{@State} copies of one fact = two truths that drift.}
  \trap{\texttt{@State} holding a plain \textbf{class}: inner mutation leaves the reference equal
        $\to$ no re-render (fine only for an \texttt{@Observable} class).}
  \trap{\texttt{if c \{ MyView() \} else \{ MyView() \}} = two identities $\to$ state resets on toggle.
        \texttt{.id(value)} changing often silently nukes state.}
  \trap{Heavy work or network calls in \texttt{body} — it runs many times; use \texttt{.task}
        (auto-cancelled on disappear) or \texttt{.onAppear}.}
\end{itemize}

\section{Timeline}
\begin{tikzpicture}[sheet]
  \tikzset{yr/.style={font=\bfseries\scriptsize, text=sheetBlue},
           ev/.style={font=\tiny, align=center, text=black!85, text width=16mm, anchor=north}}
  \draw[thick, sheetGrey, ->] (0,0) -- (8.3,0);
  \foreach \x/\y/\t in {
      0.9/iOS 13/{\texttt{@State} \texttt{@Binding}\\\texttt{@ObservedObject}\\\texttt{@EnvironmentObject}},
      3.0/iOS 14/{\texttt{@StateObject}},
      5.1/iOS 15/{\texttt{.task}\\(cancelled on\\disappear)},
      7.2/iOS 17/{\texttt{@Observable}\\\texttt{@Bindable}}}{
    \fill[sheetBlue] (\x,0) circle (0.6mm);
    \node[yr, above=1pt] at (\x,0) {\y};
    \node[ev] at (\x,-0.12) {\t};
  }
\end{tikzpicture}

\section{Remember}
\textbf{``Who wrote \texttt{= VM()}? That view writes \texttt{@StateObject}.''} State\,/\,StateObject =
\emph{mine}; Binding\,/\,Observed\,/\,Environment = \emph{borrowed}. \texttt{\$} = a pointer to the truth.


\end{multicols}

\noindent{\footnotesize\color{sheetGrey}\textit{Related:} property wrappers (\texttt{wrappedValue} /
\texttt{projectedValue}) · SwiftUI Environment \& Preferences (\texttt{PreferenceKey}, child $\to$ parent) · async work \& data flow in SwiftUI views (\texttt{.task}, restarted by \texttt{id:}) · Observation · Combine · cancellables · SwiftUI performance \& identity (\texttt{EquatableView})}

\end{document}
